\documentclass[11pt]{{ "{" }}{{cookiecutter.document_class}}{{ "}" }}

% --- Common Packages across layouts --- %
\usepackage{microtype}
\usepackage{graphicx}
\usepackage[hidelinks]{hyperref}
\usepackage{xcolor}
\usepackage{todonotes}
\usepackage{dirtytalk}
\usepackage{siunitx}

{% if cookiecutter.document_class == 'article' or cookiecutter.document_class == 'report' %}

% Article and report specific packages %

% Visual style packages

\usepackage[margin = 2.5cm]{geometry}
\usepackage{wrapfig}
\usepackage{caption}
\usepackage{subcaption}
\usepackage[nocheck]{fancyhdr}
\usepackage{titlesec}
\usepackage{blindtext}

% Maths / Science packages
\usepackage{amsfonts, amssymb, amsmath, amsthm}
\usepackage{bm}
\usepackage{physics}
\usepackage{tikz}
\usepackage{circuitikz}

% Other utility packages
\usepackage{multicol}
\usepackage{afterpage}
\usepackage{lipsum} % For testing
\usepackage[backend=biber, style=numeric]{biblatex}
\addbibresource{references.bib}

{% elif cookiecutter.document_class == 'letter' %}

% Letter specific packages %
\usepackage[margin = 3cm]{geometry}

{% endif %}

% My (John's) personal utility packages
\usepackage{packages/johns-package-setup}
\usepackage{packages/johns-preferences}

{% if cookiecutter.document_class == 'article' %}
% Article/short report template %

% Setting up headers and footers for fancyhf
\fancyhf{}
\fancyhead[R]{\today}
\fancyhead[L]{{ "{" }}{{cookiecutter.project_name}}{{ "}" }}
\fancyfoot[C]{\thepage}
\headsep = 15pt
\voffset = 10pt

% Title and introduction to the report
\title{{ "{" }}{{cookiecutter.project_name}}{{ "}" }}
\author{{ "{" }}{{cookiecutter.author_name}}\thanks{{ "{" }}{{cookiecutter.author_affiliation}}{{ "}" }}{{ "}"}}
\date{\today}

\begin{document}

\pagestyle{empty}
\newgeometry{margin = 2cm}
\pagenumbering{arabic}

\maketitle

% Set headers and footers to exist past the title page
% \afterpage{\thispagestyle{fancy}} % Was useful when not dealing with multicol
\pagestyle{fancy}

\begin{multicols}{2}

\begin{abstract}
    % Give the abstract of the document with the reason it exists, the significance of its contribution to the project and the conclusions it provides.
    \lipsum[1] 
\end{abstract}

\section{Introduction}
% Introduce the problem, and contextualise it and the objectives of the experiment/thing being reported on

% For example, remove as desired
\lipsum[0-3]

%%% DELETE SECTIONS IF INNAPROPRIATE %%%

\section{Method}
% Explain the method used in order to get the results acquired

% For example, remove as desired
\lipsum[4-7]

\section{Readings and Results}
% Give a plain, clinical explaination of the results observed, hypothesising reasons why it might be like that
\subsection{Observations}
% State any observations that were made during the experiment, and hypothesise why they may have been made

% More significant in a research context
\section{Discussion}
% Critically analyse the findings from previous sections, thereby establishing the significance of them - and in a research context then develop an argument that supports your conclusions (and thereon recommendations)

\section{Conclusions}
% Reinforce the points raised in the body sections, and ties them together to form an argument. Consolidate your assessment of your research/data and demonstrate existing gaps or problems that will be addressed in recommendations

%\section{Recommendations}
% In a research context, outline solutions to the problems highlighted in the body/conclusion of your report - also potentially suggesting what action should be undertaken, who should undertake it, what is the time frame of the action, and what costs are involved. Think SMART <---

%\section{References}
% Add the references that you have made, using data from the included bibliography.

\end{multicols}

\end{document}

{% elif cookiecutter.document_class == 'report' %}
% Report/long report template %

{% elif cookiecutter.document_class == 'letter' %}
% Letter template %



{% endif %}